%%%%%%%%%%%%%%%%%%%%%%%%%%%%%%%%%%%%%%%%%%%%%%%%%%%%%%%%%%%%%%
\section{Estruturas de Decisão, Laços e Iteradores}

%%%%%%%%%%%%%%%%%%%%%%%%%%
\begin{frame}
\frametitle{Estruturas de Decisão, Laços e Iteradores}
\begin{minipage}{0.47\textwidth}
    \begin{itemize}
        \item Definições
        \item Contexto de uso
        \item Estruturas de decisão
        \item Estruturas de repetição
        \item Iteradores
        \item Entradas e saídas
        \item Exemplos
    \end{itemize}
\end{minipage}
\begin{minipage}{0.5\textwidth}
\begin{figure}[ht!]
\begin{center}
\includegraphics[width=1.2\textwidth, height=0.40\textheight]{figures/logo_picat_alex.jpg}
\end{center}
\end{figure}
\end{minipage}
\end{frame}
%%%%%%%%%%%%%%%%%%%%%%%%%%%%%%%%%%%%%%%%%%%%%%%%%%%%%%%%%%%%%%%%%%%%%%%%%%%%%%%%%%%%%%%%%%%%%%%
%%% ATE AQUI

\begin{frame}[fragile]
\frametitle{Estruturas de Decisão, Laços e Iteradores}
    
  
  \begin{itemize}
      
      \item Ao contrário do Prolog, Picat apresenta
      conceitos e comandos da programação imperativa
      
      \pause
      \item Esta maneira  ameniza os obstáculos
      em se aprender uma linguagem com o paradigma lógico,
      tendo outros elementos conhecidos
  
          \pause
      \item Assim, Picat apresenta estruturas clássicas
      como: 
      \begin{itemize}
        \item  \texttt{if--then--end}, \texttt{if--then--else--end}, \texttt{if--then--elseif--then--....end}
        \item  \texttt{foreach}
        \item  \texttt{while}
        \item  \texttt{do--while}
        \item Bem como a atribuição, `\texttt{:=}', já discutida
      \end{itemize}
   \end{itemize}

       
\end{frame}    
     
    %%%%%%%%%%%%%%%%
    
\begin{frame}[fragile]
\frametitle{Estrutura Condicionais}
    
    \begin{itemize}
        
        \item Picat implementa uma  estrutura condicional explícita (na programação em lógica, voce faz isto implicitamente)
        
        \pause
        \item Sua notação é:\\
        
        \begin{tabbing}
            aa \= aaa \= aaa \= aaa \= aaa \= aaa \= aaa \kill
            \> \texttt{if (Exp) then}\\
            \> \> \texttt{Ações} \\
            \> \texttt{else}\\
            \> \> \texttt{Ações}\\
            \> \vdots\\
            \> \texttt{end} 
        \end{tabbing}
        
        \pause
        \item Onde \emph{Exp} é uma expressão lógica  avaliada como verdadeira ou  falsa.
     
     \pause   
        \item A última ação antes de um \emph{else} ou \emph{end} não deve ter  vírgula nem ponto e vírgula ao final da linha.
  
     \pause   
        \item Tem-se ainda o \texttt{elseif} que pode estar embutido no comando \texttt{if--then--else--end}
        
    \end{itemize}
\end{frame}    
    
    
\begin{frame}[fragile]

\begin{block}{Exemplo: \texttt{if--then--else--end} }
     
\begin{lstlisting}[frame=single]
    ,
if (X <= 100) then
        println("X e menor que 100")
    elseif (X <= 1000 && X >= 500) then
        println("X estah entre 500 e 1000")
    else
        println("X estah abaixo de  500")
    end
    ,
\end{lstlisting}
        
\end{block}
\end{frame}    

    
    
%%%%%%%%%%%%%%%%%%%%%%%%%%%%%%%%%%%%%%%%%%%%%%%%%%%%%%%%%%%%%%%    
       
    
\begin{frame}[fragile]
\frametitle{Estruturas de Repetições}
       

\begin{itemize}
    
    \item Picat também implementa 3 estruturas de repetição:
    \texttt{foreach}, \texttt{while} e \texttt{do-while}
    
    \pause
    \item O laço do  \texttt{foreach}  \textbf{itera} sobre termos simples e compostos
    
    \pause
    \item O  \texttt{while} repete um conjunto de ações enquanto uma 
     condição for verdadeira.
     
        \pause
    \item A condição pode ser simples ou combinada: ver exemplos
     
     \pause
    \item O laço \texttt{do-while} é análogo ao  \texttt{while}, porém ele 
    sempre executa pelo menos uma vez
\end{itemize}
\end{frame}    


%%%%%%%%%%%%%%%%%%%%%%%%%%%%%%%%%%%%%%%%%%%%%%%%%%%%%%%%%%%%%%%    
\begin{frame}[fragile]
\frametitle{Estruturas de Repetições: \texttt{foreach}}

    \begin{itemize}
        \item Um laço \texttt{foreach} tem a seguinte forma:
        
        \begin{tabbing}
            aa \= aaa \= aaa \= aaa \= aaa \= aaa \= aaa \kill
            \> \texttt{foreach ($E_1$ in $D_1$, $Cond_1$, $\ldots$, $E_n$ in $D_n$, $Cond_n$)}  \\
            \> \> $Metas$ \\
            \>  \texttt{end} 
        \end{tabbing}
    
   \pause    
   Esta notação é dada por:        
    \item  $E_i$ é um \emph{padrão de iteração} ou \emph{iterador}. 
        
    \item  $D_i$ é  uma expressão de \emph{valor composto}. Exemplo: uma lista de valores
        
     \item  $Cond_i$ é uma condição opcional sobre os \textbf{iteradores} $E_1$ até $E_i$.
        
     \item O \texttt{foreach} pode conter múltiplos iteradores usando o ``\texttt{in}'' 
     Caso isso ocorra, o compilador interpreta isso como diversos laços aninhados.
        
    \end{itemize}
\end{frame}    


\begin{frame}[fragile]

\begin{block}{Exemplo: \texttt{foreach} }
     
\begin{lstlisting}[frame=single]
 imp_tracejados (N) =>
   nl,
   foreach(I in 1..N)
     printf("=")
   end,
   nl.   
..................................................................  
Picat> cl('backtracking_ex_02').
Compiling:: backtracking_ex_02.pi
** Warning: singleton variables (backtracking_ex_02.pi, 50-55): I
backtracking_ex_02.pi compiled in 3 milliseconds
loading...

yes

Picat> imp_tracejados(30).

==============================
\end{lstlisting}
        
\end{block}

\textcolor{red}{O \texttt{I} é iterado com valores de \texttt{1} a \texttt{N}}
\end{frame}    


%%%%%%%%%%%%%%%%%%%%%%%%%%%%%%%%%%%%%%%%%%%%%%%%%%%%%%%%%%%%%%%    
\begin{frame}[fragile]
\frametitle{Estruturas de Repetições: \texttt{while}}


\begin{itemize}
        \item O laço do \texttt{while} tem a seguinte forma:
        
        \begin{tabbing}
            aa \= aaa \= aaa \= aaa \= aaa \= aaa \= aaa \kill
            \> \texttt{while ($Cond$)} \\
            \> \> $Metas$  \\
            \>  \texttt{end}
        \end{tabbing} 
        
        \item Enquanto a expressão lógica \emph{Cond} for verdadeira, o conjunto de \texttt{Metas} é executado.
        
\end{itemize}
\end{frame}    


\begin{frame}[fragile]

\begin{block}{Exemplo: \texttt{while} }
     
\begin{lstlisting}[frame=single]
  
laco_02 =>  
    I = 1,  
    while (I <= 9)  
        println(I),  
        I := I + 2  
    end.  
\end{lstlisting}
        
\end{block}
\end{frame}    

%%%%%%%%%%%%%%%%%%%%%%%%%%%%%%%%%%%%%%%%%%%%%%%%%%%%%%%%%%%%%%%    
\begin{frame}[fragile]
\frametitle{Estruturas de Repetições: \texttt{do-while}}

    \begin{itemize}
        
        \item O laço \texttt{do-while} tem a seguinte forma:
        
        \begin{tabbing}
            aa \= aaa \= aaa \= aaa \= aaa \= aaa \= aaa \kill
            \> \texttt{do} \\
            \> \> $Metas$  \\
            \>  \texttt{while ($Cond$)}
        \end{tabbing} 
        
        \item Ao contrário do  \texttt{while} o iterador \texttt{do-while}
        vai executar \texttt{Metas} pelo menos uma vez antes de avaliar \texttt{Cond}.
        
    \end{itemize}
    
\end{frame}


\begin{frame}[fragile]

\begin{block}{Exemplo: \texttt{do-while} }

\begin{lstlisting}[frame=single]
     
laco_03 =>  
     J = 6,  
    do  
        println(J),  
        J := J + 1  
    while (J <= 5).  
\end{lstlisting}
        
\end{block}
\end{frame}    


%%%%%%%%%%%%%%%%%%%%%%%%%%%%%%%%%%%%%%%%%%%%%%%%%%%%%%%%%%%%%%%%%%%%%

\subsection{Funções e Predicados Especiais}


\begin{frame}[fragile]
%[c,allowframebreaks]

\frametitle{Funções e Predicados Especiais}
    % Compreensão de lista
    % Entrada e Saida de Dados
\begin{itemize}
    
    \item Há algumas funções e predicados especiais em Picat 
    que necessitam de algum cuidado:
    
    \pause
    \begin{itemize}
      \item \textcolor{magenta}{Geração (compreensão) de listas e vetores}
      \item  \textcolor{magenta}{Entradas e saídas de dados}
    \end{itemize}
    
    \pause
    \item Na verdade, já fizemos uso delas, porém sem a ênfase de que
    são funções ora predicados.
    
    \pause
    \item \textcolor{magenta}{Um \textbf{atalho} consider\'avel na programação com Picat!}
\end{itemize}
\end{frame}      


%%%%%%%%%%%%%%%%%%%%%%%%%%%%%%%%%%%%%%%%%%%%%%%%
\begin{frame}[c,allowframebreaks]

\frametitle{Compreensão de Listas e Vetores}    
    
\begin{itemize}    
    \item A função de \underline{\textit{compreensão de listas e vetores}} é uma 
    função especial que  permite a fácil criação de listas ou vetores.
    
    \item Sua notação é:
    
    \begin{tabbing}
        aa \= aaa \= aaa \= aaa \= aaa \= aaa \= aaa \kill
        \> \> \texttt{[$T$ : $E_1$ \texttt{in} $D_1$, $Cond_1$, $\ldots$, $E_n$ in $D_n$, $Cond_n$]} 
    \end{tabbing}
    
    \item Onde, $T$ é uma expressão  adicionada a lista, cada $E_i$ é um 
    iterador sobre $D_i$, o qual é um termo ou expressão, 
     e $Cond_i$ é uma condição sobre cada iterador de $E_1$ até $E_i$.
    
    \item Há uma seção dedicada a listas. Voltaremos ao assunto.
    
    \framebreak

    \item Esta função pode gerar um vetor também, a notação é um pouco diferente:
    
    \begin{tabbing}
        aa \= aaa \= aaa \= aaa \= aaa \= aaa \= aaa \kill
        \> \> \texttt{\{$T$ : $E_1$ \texttt{in} $D_1$, $Cond_1$, $\ldots$, $E_n$ in $D_n$, $Cond_n$\}} 
    \end{tabbing}
    
    \item Neste caso, os delimitadores são $\{$ e $\}$ de um vetor

\end{itemize}

\end{frame}      

% %%%%%%%%%%%%%%%%%%%%%%%%%%%%%%%%%%%%%%%%%    
 \begin{frame}[fragile]
 \frametitle{Compreensão de Listas e Vetores: Exemplo}
\begin{footnotesize}

\begin{verbatim}
main =>
 L = [(A, I) : A in [a, b], I in 1 .. 2],
 V = {(I, A) : A in [a, b], I in 1 .. 2},
 printf("\nLista: %w  \nVetor: %w\n", L, V),
 imp_vetor(V),
 imp_vetor(L).

imp_vetor (M) => 
 Tam = M.length,  %% tamanho de M
  nl,
   foreach(I in 1  .. Tam )
     printf("V_ou_L(%d):%w \t" , I, M[I] )
    end,
  nl.
%%%% $picat vetor_exemplo_01.pi
\end{verbatim}
\end{footnotesize}
 
\end{frame}

% %%%%%%%%%%%%%%%%%%%%%%%%%%%%%%%%%%%%%%%%%    
 \begin{frame}[fragile]
 \frametitle{Saída -- Compreensão de Listas e Vetores}
\begin{footnotesize}
\begin{verbatim}
$ picat vetor_exemplo_01.pi 

Lista: [(a,1),(a,2),(b,1),(b,2)]  
Vetor: {(1,a),(2,a),(1,b),(2,b)}

V_ou_L(1):1,a 	V_ou_L(2):2,a 	V_ou_L(3):1,b 	V_ou_L(4):2,b 	

V_ou_L(1):a,1 	V_ou_L(2):a,2 	V_ou_L(3):b,1 	V_ou_L(4):b,2 	
[hulk@SAGUACU picat]$ 
\end{verbatim}
\end{footnotesize}
 
\end{frame}

%%%%%%%%%%%%%%%%%%%%%%%%%%%%%%%%%%%%%%%%%%%%%%%%%%%%%%
\begin{frame}[c,allowframebreaks]

\frametitle{Leitura e Escrita}    

\begin{itemize}
        
\item Picat tem diversas  funções de leitura de valores,
        que serve tanto para ler de uma console \texttt{stdin},
        como de um arquivo qualquer.
        
\item Aos usuários de Prolog, aqui não precisamos do delimitador
final de `\textbf{.}'   ao final de uma leitura. 

\item Válido quando editamos no interpretador, o `\textbf{.}' final é opcional
        
\framebreak        
\item As mais importantes são:

\begin{itemize}
  
  \item \texttt{read\_int($FD$) = $Int$} $\Rightarrow$ Lê um \textit{Int} do 
  arquivo $FD$.
  
  \item \texttt{read\_real($FD$) = $Real$} $\Rightarrow$ Lê um \textit{Float} do 
  arquivo $FD$.
  
  \item \texttt{read\_char($FD$) = $Char$} $\Rightarrow$ Lê um \textit{Char} do 
  arquivo $FD$.
  
  \item \texttt{read\_line($FD$) = $String$} $\Rightarrow$ Lê uma \textit{Linha} 
  do arquivo $FD$.
  
\end{itemize}

\item Caso se deseja ler da console, padrão \texttt{stdin}, $FD$, o nome do descritor de arquivo,  pode ser omitido.

\framebreak

\item Os dois predicados mais importantes para saída de dados, são 
\texttt{write} e \texttt{print}.

\item Cada um destes predicados tem três variantes, são eles:

\begin{itemize}

\item \texttt{write($FD$, $T$)} $\Rightarrow$ Escreve um termo $T$ no arquivo 
$FD$.

\item \texttt{writeln($FD$, $T$)} $\Rightarrow$ Escreve um termo $T$ no arquivo 
$FD$, e pula uma linha ao final do termo.

\item \texttt{writef($FD$, $F$, $A\ldots$)} $\Rightarrow$ Este predicado é usado 
para escrita formatada para um arquivo $FD$, onde $F$ indica uma série de 
formatos para cada termo contido no argumento $A\ldots$. O número de argumentos 
não pode exceder 10.

\end{itemize}
        
\framebreak
        
\item Analogamente, para o predicado \texttt{print}, temos:

\begin{itemize}

  \item \texttt{print($FD$, $T$)} $\Rightarrow$ Escreve um termo $T$ no arquivo 
  $FD$.
  
  \item \texttt{println($FD$, $T$)} $\Rightarrow$ Escreve um termo $T$ no arquivo 
  $FD$, e pula uma linha ao final do termo.
  
  \item \texttt{printf($FD$, $F$, $A\ldots$)} $\Rightarrow$ Este predicado é usado 
  para escrita formatada para um arquivo $FD$, onde $F$ indica uma série de 
  formatos para cada termo contido no argumento $A\ldots$. O número de argumentos 
  não pode exceder 10.

\end{itemize}

\item Caso queira escrever para \texttt{stdout}, o nome do  $FD$, pode ser omitido.

\end{itemize}
    
\end{frame}

%%%%%%%%%%%%%%%%%%%%%%%%%%%%%%%%%%%%%%%%%%%%%%%%%%%%%%%%%%%%%%%%%%%%%

\begin{frame}
{Tabela de Formatação para Escrita}

Apenas os mais importantes, há outros como: hexadecimal, notação científica,
etc. Ver no apêndice do Guia do Usuário.    
\begin{table}[h]
    \begin{tabular}{l|l}
        \hline  \hline
        \textbf{Especificador} & \textbf{Saída} \\
        \hline 
        \hline 
        \texttt{\%\%} & Sinal de Porcentagem \\
        \texttt{\%c} & Caráctere \\
        \texttt{\%d \%i} & Número Inteiro Com Sinal \\
        \texttt{\%f} & Número Real \\
        \texttt{\%n} & Nova Linha \\
        \texttt{\%s} & \textit{String} \\
        \texttt{\%u} & Número Inteiro Sem Sinal \\
        \texttt{\%w} & Termo \textbf{qualquer} \\
        \hline   \hline
    \end{tabular}
\end{table}
\end{frame}

%%%%%%%%%%%%%%%%%%%%%%%%%%%%%%%%%%%%%%%%%%%%%%%%%%%%%%%%%%%%%%%%%%%%%

\begin{frame}

\frametitle{Comparação entre \texttt{write} e \texttt{print}}
     
\begin{table}[h]
    \centering
    \begin{tabular}{ l|l|l|l }
        \hline         \hline
        Dados $\Rightarrow$  & "abc" & [a,b,c] & 'a@b'\\
        \hline         \hline
        \texttt{write}   & [a,b,c] & [a,b,c] & 'a@b' \\
        \hline
        \texttt{writef}  & [a,b,c] (\%s) & abc (\%w) & 'a@b' (\%w) \\
        \hline
        \texttt{print}   & abc & abc & a@b \\
        \hline
        \texttt{printf}  & abc (\%s) & abc (\%w) & a@b (\%w) \\ 
        \hline \hline
    \end{tabular}
\end{table}

\end{frame}

%%%%%%%%%%%%%%%%%%%%%%%%%%%%%%%%%%%%%%%%%%%%%%%%%%%%%%%%%%%%%%%%%%%%%

\begin{frame}[fragile]

\frametitle{Exemplos}
    
    \begin{block}{Condicionais}
        \begin{lstlisting}[frame=single]
main =>
    X = read_int(),
    if(X <= 100)then
        println("X e menor que 100")
    else
        println("X nao e menor que 100")
    end.
 \end{lstlisting}
    
    \end{block}
\end{frame}

\begin{frame}[fragile]

\frametitle{Exemplos -- Repetições}
    
   \begin{lstlisting}[frame=single]
main =>
    X = read_int(),
    println(x=X),
    while(X != 0)
        X := X - 1,
        println(x=X)
    end
.
   \end{lstlisting}
        
   \begin{lstlisting}[frame=single]
main =>
    X = read_int(),
    Y = X..X*3,
    foreach(A in Y)
        println(A)
    end.
        \end{lstlisting}
    
\end{frame}

\begin{frame}[fragile]
\frametitle{Exemplos -- Construindo  Listas e Vetores}    
       
       
\begin{footnotesize}
\begin{verbatim}
import util. % use split
main => 
        le_vetor_01(X1),
        printf("\nVETOR LIDO: %w ", X1),
        le_vetor_02(X2),
        printf("\nVETOR LIDO: %w ", X2),
        le_lista_01(Y),
        printf("\nLISTA LIDA: %w ", Y),
        le_lista_02(W),
        printf("\nLISTA LIDA 2: %w " , W) .
\end{verbatim}
\end{footnotesize}       

\begin{small}
\begin{itemize}
  \item \textcolor{red}{Este exemplo reúne muitos conceitos desta seção.}  
  \item \textcolor{red}{\url{https://github.com/claudiosa/CCS/blob/master/picat/input_output_exemplos/leitura_vetores_listas.pi}}  
\end{itemize}
\end{small}

\end{frame}


\begin{frame}[fragile]
\frametitle{Exemplo -- Leitura de Vetores}    

\begin{footnotesize}
\begin{verbatim}
le_vetor_01 ( V ) =>
  printf("\nDIGITE tamanho da entrada: "),
  Tam = read_int(),
  V = new_array( Tam ), % cria um vetor
  printf("\nDIGITE os %d VALORES do vetor:", Tam),
  foreach (I in 1..Tam)
        V[I] = read_int()
  end,
  printf("\nVETOR: %w ", V).
  
le_vetor_02 ( V ) =>  
  printf("\nLendo um vetor qualquer de inteiros na linha: "),
  V = { to_int(W) : W in read_line().split() }.
  % OU
  %L = [ to_int(W) : W in read_line().split()],
  %V = to_array( L ).
\end{verbatim}
\end{footnotesize}       
\end{frame}


\begin{frame}[fragile]
\frametitle{Exemplo -- Leitura de  Listas}    

\begin{footnotesize}
\begin{verbatim}
le_lista_01 ( L ) =>  
  printf("\nLendo lista de inteiros na linha: "),
  L = [ to_int(W) : W in read_line().split()].
  
le_lista_02 (List) =>  
    printf("\nLista inteiros e 0 encerra: "),
    L := [] ,
    E := read_int() ,
    while (E != 0)  
        L := [E|L],  
        E := read_int()  
    end,  
    List = L .
\end{verbatim}
\end{footnotesize}       
\textcolor{red}{Volte neste exemplo após a seção de Listas.}          
\end{frame}


\begin{frame}[fragile]
\frametitle{Saída -- via \texttt{main}}    

\begin{footnotesize}
\begin{verbatim}
$ picat leitura_vetores_listas.pi 
DIGITE tamanho da entrada: 3
DIGITE os 3 VALORES do vetor: 3 4 5
VETOR LIDO: {3,4,5} 
Lendo um vetor qualquer de inteiros na linha: 9 8 7 6 
VETOR LIDO: {9,8,7,6} 
Lendo lista de inteiros na linha: 1 2 3 4 
LISTA LIDA: [1,2,3,4] 
Lista inteiros e 0 encerra: 1 2 3 7 0 1 2 3 5
LISTA LIDA 2: [7,3,2,1] 
.... removi algumas linhas em branco
\end{verbatim}
\end{footnotesize}       
\end{frame}

%%%%%%%%%%%%%%%%%%%%%%%%%%%%%%%%%%%%%%%%%%%%%%%%%%%%%%%%%%%%%%%%%%%%%

\begin{frame}[fragile]
\frametitle{Reflexões}

\begin{itemize}

  \item Esta seção avança na sintaxe do Picat e precisa 
  ser praticada
  
   \pause
  \item Como este conteúdo se assemelha as LPs clássicas, como exercício,
  voce está apto a fazer alguns algoritmos de outras linguagens.

  \pause 
  \item Se antes o Prolog era complicado para \textit{entradas} e \textit{saídas}, com Picat,
  tudo ficou semelhante as demais linguagens imperativas


  \pause
  \item \textcolor{red}{\textbf{Cuidado}: as funções de \textit{entradas} e \textit{saídas}, e outras do sistema como
  \textit{time}, etc,  \textbf{nunca} falham e nem aceitam \textit{backtracking} 
  definidas para elas. Em caso de falha do predicado que as contém, estas são \textbf{silenciosas},
  apresentando um simples \textbf{no} ou \textbf{false}!}
  
  \pause
  \item \textcolor{green}{Felizmente os erros e problemas de sintaxe  são prontamente reportados}
  
  \pause
  \item Em \url{https://github.com/claudiosa/CCS/tree/master/picat/input_output_exemplos}
  tem vários exemplos avançados de entradas e saídas 
  
  \pause
  \item Mãos à obra!
 
\end{itemize}

\end{frame}
